% P05.1 one-dimensional regularity convention repair proof, 25 September 2026.

\begin{proof}[\Cref{thm:frontier-polyhedral-one-dimensional-real-regularity}]
Put $a:=\gamma+u$ and $P:=P_\gamma(u)$.  If $C\geq-\inf_{\mathbf X}u$,
then $\gamma-C$ is an admissible competitor, so $P$ is finite.  By
\cite[Proposition~6.3]{BMPS26}, it is polyhedrally convex, satisfies
$P\leq a$, and differs from $\gamma$ by a bounded function.  Insert one
rational vertex in every component of $\mathbf X$ having no vertex; this does
not change the assertion.

For a bounded edge $e=[v_0,v_1]$ of $\Pi$, let $L_e$ be the affine function
on $e$ with values $P(v_0)$ and $P(v_1)$ at its endpoints.  Convexity gives
$P|_e\leq L_e$.  Since $a|_e$ is affine and $P(v_i)\leq a(v_i)$ for
$i=0,1$, one also has $L_e\leq a|_e$.

Now let $e$ be an unbounded edge with endpoint $v$, and use its primitive
integral-affine coordinate $t\geq0$.  The function
\[
                    F_e(t):=P(t)-\gamma(t)
\]
is bounded and convex.  It is nonincreasing: if a right derivative were
positive at some point, monotonicity of the one-sided derivatives of a convex
function would make $F_e$ unbounded above.  Hence
\[
 P|_e\leq L_e:=\gamma|_e+P(v)-\gamma(v).
\]
The bounded affine function $u|_e$ is constant.  Evaluating $P\leq a$ at
$v$ therefore gives $P(v)-\gamma(v)\leq u|_e$, and hence $L_e\leq a|_e$.

Define $Q$ by $Q|_e=L_e$ on every edge and $Q=P$ at every vertex.  The
definitions agree at the endpoints, so $Q$ is continuous.  It is real
piecewise affine on $\Pi$, satisfies
\[
                         P\leq Q\leq a,
 \qquad Q-\gamma=O(1),
\]
and is polyhedrally convex.  Indeed, a polyhedrally convex combination whose
centre lies in the relative interior of an edge uses only that edge, where
$Q$ is affine.  If its centre is a vertex $v$, then $Q(v)=P(v)$ and
$Q\geq P$ at every point entering the combination, so the Jensen inequality
for $P$ implies the one for $Q$.

Thus $Q$ is an admissible competitor for $P_\gamma(u)$.  The definition of
the envelope gives $Q\leq P$, while the construction gives $P\leq Q$.
Therefore $P=Q$, proving that $P$ is real piecewise affine on $\Pi$.  Since
it is polyhedrally convex and has $\gamma+O(1)$ growth, it belongs to
$\mathrm{PAPC}_{\mathbf R}(\mathbf X,\gamma)$ and hence to
$\mathrm{PC}_{\mathrm{reg},\mathbf R}(\mathbf X,\gamma)$.
\end{proof}

\begin{proof}[\Cref{prop:frontier-polyhedral-integral-regularity-counterexample}]
The displayed $u$ is bounded and has integral slopes on the rational
subdivision with vertices $0,1,2$.  Put $a:=\gamma+u$.  The displayed
candidate $P$ is convex, has $\gamma+O(1)$ growth, and lies below $a$.
Conversely, every convex competitor $q\leq a$ satisfies $q(0)\leq0$ and
$q(2)\leq1$.  Convexity therefore gives $q(x)\leq x/2$ on $[0,2]$, while
the obstacle inequality gives $q\leq P$ outside that interval.  Hence the
candidate is $P_\gamma(u)$.

Let $\psi\in\mathrm{PAPC}_{\mathbf Z}(\mathbf R,\gamma)$.  Its restriction
to $[0,2]$ is convex and piecewise affine with integral slopes.  The left
derivative $m:=\psi'_{-}(1)$ is therefore an integer.  If $m\leq0$, convexity
gives
\[
                         \psi(1)-\psi(0)\leq0.
\]
Writing $e_j:=\psi(j)-P(j)$, this implies
\[
                         e_1-e_0\leq-\frac12,
\]
so $\max\{|e_0|,|e_1|\}\geq1/4$.  If $m\geq1$, monotonicity of one-sided
derivatives gives
\[
                         \psi(2)-\psi(1)\geq1,
\]
and hence $e_2-e_1\geq1/2$, so
$\max\{|e_1|,|e_2|\}\geq1/4$.  These two cases exhaust the possible integral
values of $m$, and therefore
\[
             \lVert\psi-P\rVert_{L^\infty(\mathbf R)}\geq\frac14.
\]

No sequence in the literal integral-slope
$\mathrm{PAPC}_{\mathbf Z}(\mathbf R,\gamma)$ can converge uniformly to
$P$.  In particular $P$ does not belong to the corresponding regularizable
class.  The bounded literal piecewise-affine obstacle $u$ therefore violates
condition~(1) of \cite[Theorem~6.4]{BMPS26}.
\end{proof}
